\title{Controlling Text-to-Image Diffusion by Orthogonal Finetuning}

\begin{abstract}
Recent advancements in large text-to-image diffusion models have enabled the synthesis of photorealistic images conditioned on textual descriptions. A pressing research question is how to effectively steer or customize these robust models for a variety of downstream tasks. In this context, we present Orthogonal Finetuning (OFT), a methodologically grounded finetuning strategy designed for the adaptation of text-to-image diffusion models. Distinct from previous approaches, OFT is theoretically shown to conserve hyperspherical energy, which reflects the mutual relationship between neurons mapped onto a unit hypersphere. Maintaining this property is found to be vital for sustaining the semantic generative prowess of these models. To further enhance stability during finetuning, we introduce Constrained Orthogonal Finetuning (COFT), which incorporates a supplementary radius constraint within the hyperspherical manifold. Our investigation centers on two principal finetuning scenarios in text-to-image generation: subject-driven generation, where the objective is to synthesize images specific to a given subject using a limited set of images and a prompt; and controllable generation, in which models are adapted to ingest auxiliary control signals. Through extensive empirical evaluation, our OFT paradigm consistently yields superior results over established baselines in terms of both generation quality and training efficiency.
\end{abstract}

\section{Introduction}

State-of-the-art text-to-image diffusion frameworks~\cite{saharia2022photorealistic,ramesh2022hierarchical,rombach2022high} have demonstrated remarkable competency in leveraging textual input for photorealistic image creation. However, while their results are visually compelling, the interpretability and granularity of text conditioning often fall short for tasks demanding precise and detailed control over generated content. This work focuses on addressing these limitations by concentrating on two specialized categories of text-to-image generation tasks:

\begin{itemize}[leftmargin=*,nosep]

    \item \textbf{Subject-driven generation}~\cite{ruiz2023dreambooth}: This task entails synthesizing novel images of a particular subject in various scenarios, using only a handful of example images of the subject in combination with a textual prompt.
    \item \textbf{Controllable generation}~\cite{zhang2023adding,mou2023t2i}: In this setting, the model is provided with additional modalities of guidance (such as canny edge detections or segmentation annotations) to influence the output, alongside the standard text prompt.
\end{itemize}

Both tasks fundamentally address the challenge of how to adapt text-to-image diffusion models to new tasks while maintaining their pretrained generative capabilities. The key criteria for a robust finetuning approach include: (1) \emph{training efficiency}, characterized by a minimal set of trainable parameters and reduced training epochs, and (2) \emph{generalizability preservation}, involving the retention of high-quality and diverse generative outputs. Conventional finetuning is typically achieved either by directly adjusting existing neuron weights with a low learning rate (\emph{e.g.}, \cite{ruiz2023dreambooth}) or by appending lightweight, re-parameterized modules to the original model (\emph{e.g.}, \cite{hulora2022,zhang2023adding}). While these methods are straightforward, they do not offer explicit guarantees of maintaining generative performance from pretraining. Furthermore, there remains a gap in systematic approaches for crafting effective finetuning techniques and tuning hyperparameters such as training duration. One principal obstacle is the absence of a clear metric for measuring the extent to which pretrained generative abilities are preserved. Most current methods implicitly equate a smaller Euclidean distance between the parameters of the finetuned and pretrained models with better preservation of original abilities, motivating the frequent use of very low learning rates. Although this assumption may sometimes hold, we contend that Euclidean proximity alone is insufficient to represent semantic preservation. Thus, introducing a more structurally-informed metric for assessing the divergence between finetuned and original models could substantially enhance both the conservation of pretraining performance and the reliability of the finetuning process.

Drawing upon empirical findings that hyperspherical similarity effectively encodes semantic relationships~\cite{liu2017deep,liu2018decoupled,chen2020angular}, our approach employs hyperspherical energy~\cite{liu2018learning} as a means to capture the pairwise relational patterns among neurons. Hyperspherical energy is calculated by summing the hyperspherical similarity (\emph{e.g.}, cosine similarity) across all neuron pairs within a given layer, which quantifies the degree to which neurons are uniformly distributed over the unit hypersphere~\cite{liu2021learning}. Our central hypothesis posits that an optimal finetuned model should exhibit only a negligible change in hyperspherical energy relative to its pretrained counterpart. While one might consider simply introducing a regularization term to enforce invariant hyperspherical energy during finetuning, such an approach does not guarantee the minimization of its variation. Instead, we exploit a unique invariance characteristic of hyperspherical energy—specifically, the pairwise hyperspherical similarity is rigorously maintained when a common orthogonal transformation is applied to all neurons. Building upon this property, we introduce Orthogonal Finetuning~(OFT), an adaptation technique for large text-to-image diffusion models that maintains hyperspherical energy during downstream task adaptation. The main concept involves learning a single orthogonal transformation per layer for the neurons so as to preserve all pairwise angles among them. This framework can also be interpreted as rotating the underlying coordinate system for neurons within a layer. Additionally, recognizing that reducing the Euclidean distance between the parameters of the finetuned and pretrained models favors the retention of pretraining knowledge, we extend OFT to a variant named Constrained Orthogonal Finetuning~(COFT). COFT enforces the finetuned parameters to reside within a fixed-radius hypersphere centered at the pretrained neuron positions.

The rationale behind employing orthogonal transformations in finetuning neurons is partly inspired by insights from the 2D Fourier transform, where images are represented in terms of magnitude and phase spectra. The phase spectrum, encoding the angular relationship between inputs and bases, captures most of the semantic content. Notably, one can reconstruct an image by combining its phase spectrum with a randomly chosen magnitude spectrum, while still maintaining its semantic properties. This observation indicates that adjusting the orientation of neurons is central to modifying the generated image semantically—an essential aspect of both subject-driven and controlled image generation. However, excessive modification of neuron directions can lead to degradation of the model's pretrained generative capacity due to the increased degrees of freedom. To mitigate this, we propose maintaining the angular relationships among all neuron pairs, premised on the belief that these inter-neuron angles encode critical knowledge within neural networks. Accordingly, we advocate for learning an orthogonal transformation shared across layers to ensure that hyperspherical energy remains constant within each layer.

Further conceptual motivation is drawn from orthogonal over-parameterized training~\cite{liu2021orthogonal}, which involves training classification networks from the ground up by applying orthogonal transformations to networks initialized randomly. Such random initialization is known to yield, on average, a small hyperspherical energy, and the primary objective of~\cite{liu2021orthogonal} is to maintain this low energy throughout training (as lower hyperspherical energy correlates with enhanced generalization in classification tasks~\cite{liu2018learning,lin2020regularizing}). The findings of~\cite{liu2021orthogonal} demonstrate that orthogonal transformations offer ample flexibility for training robust classification networks. Our work diverges by concentrating on finetuning text-to-image diffusion models to achieve superior control and stronger generative outcomes in downstream tasks. The contrast between OFT and~\cite{liu2021orthogonal} manifests in two principal ways. Firstly, while~\cite{liu2021orthogonal} is formulated to minimize hyperspherical energy, OFT seeks to retain the same hyperspherical energy as the pretrained model, ensuring that finetuning does not compromise the original model structure. For diffusion models, pursuing minimal hyperspherical energy may in fact disrupt the semantic organization of representations. Secondly, OFT introduces constraints to reduce deviation from the pretrained model, giving rise to its constrained formulation—such constraints are absent in~\cite{liu2021orthogonal}. Achieving an optimal balance between model adaptability and structural fidelity is critical in finetuning, and we posit that our OFT framework accomplishes this balance effectively. Our main contributions are summarized below:

\begin{itemize}[leftmargin=*,nosep]

    \item We introduce a new finetuning technique—Orthogonal Finetuning (OFT)—designed to enhance the controllability of text-to-image diffusion models. To promote further stability during adaptation, we also develop a constrained extension that restricts the angular displacement relative to the original pretrained network.
    \item Unlike traditional finetuning strategies, OFT refines models while rigorously conserving the hyperspherical energy—a property we observe empirically to be closely associated with the preservation of generative semantics from the original model.
    \item We utilize OFT for two central applications: subject-driven generation and controllable generation. Through extensive experimentation, we illustrate that our method surpasses existing approaches in several aspects, including image quality, convergence rate, and finetuning robustness. Additionally, OFT demonstrates superior data efficiency, achieving effective convergence with notably fewer training examples and iterations.
    \item For the controllable generation task, we present a novel control consistency metric to quantify controllability. This metric operates by extracting the control cue from synthesized outputs and comparing it to the corresponding input control information, thereby providing further evidence for OFT's enhanced controllability.
\end{itemize}

\section{Related Work}

\textbf{Text-to-image diffusion models}. Recent advances~\cite{saharia2022photorealistic,ramesh2022hierarchical,rombach2022high,gu2022vector,nichol2022glide} have significantly improved text-to-image synthesis, fueled by developments in diffusion-based generative modeling~\cite{ho2020denoising,song2021score,song2021denoising,dhariwal2021diffusion} as well as in vision-language joint representations~\cite{su2020vlbert,lu2019vilbert,radford2021learning,wang2021simvlm,li2022blip,li2023blip,alayrac2022flamingo,schuhmann2022laion}. Models such as GLIDE~\cite{nichol2022glide} and Imagen~\cite{saharia2022photorealistic} operate by training diffusion mechanisms within the pixel space. GLIDE jointly optimizes the text encoder with a diffusion prior based on paired image-text data, whereas Imagen employs a pretrained, frozen text encoder. In contrast, Stable Diffusion~\cite{rombach2022high} and DALL-E2~\cite{ramesh2022hierarchical} conduct training in the latent space, with Stable Diffusion utilizing VQ-GAN~\cite{esser2021taming} to build a latent visual codebook, and DALL-E2 leveraging CLIP~\cite{radford2021learning} for constructing a shared latent embedding between images and textual descriptions. Besides diffusion-based architectures, work on generative adversarial networks~\cite{reed2016generative,zhang2017stackgan,xu2018attngan,li2019controllable} and autoregressive frameworks~\cite{ramesh2021zero,ding2021cogview,wu2022nuwa,yuscaling} has further expanded text-to-image synthesis research. The proposed OFT is a flexible, model-agnostic adaptation strategy applicable to any text-to-image diffusion architecture.

\textbf{Subject-driven generation}. Approaches to mitigate unwanted subject modification include augmenting models with user-provided masks as an additional conditioning input~\cite{avrahami2022blended,nichol2022glide}. Inversion approaches~\cite{choi2021ilvr,dhariwal2021diffusion,ramesh2022hierarchical,gal2022image} adjust contextual attributes without altering the main subject. Techniques such as~\cite{hertz2022prompt} can facilitate both localized and holistic edits without relying on input masks. However, these strategies often fall short in accurately preserving distinctive subject identities. Pivotal Tuning~\cite{roich2022pivotal} locally adapts a generator near an initial latent inversion, incorporating an additional constraint to maintain identity features. Similarly,~\cite{nitzan2022mystyle} customizes a facial generative prior using a dataset of an individual's facial images, while~\cite{casanova2021instance} supports variation synthesis for specific entities, occasionally at the expense of fine instance detail. DreamBooth~\cite{ruiz2023dreambooth} achieves personalized synthesis by finetuning with reconstruction and class-aware prior preservation losses, utilizing a unique token and a limited set of subject images. The OFT method adopts the overall DreamBooth pipeline, but diverges by employing orthogonal parameter updates instead of conventional low-rate finetuning, facilitating more stable and precise model adaptation.

\textbf{Controllable generation}. Image-to-image translation is fundamentally a task involving controllable generation, where the generation process is guided by input signals. Earlier approaches have predominantly relied on conditional generative adversarial networks~\cite{isola2017image,zhu2017toward,wang2018high,park2019semantic,choi2018stargan}. More recently, diffusion models have been introduced for this purpose as well~\cite{saharia2022palette,voynov2022sketch,wang2022pretraining}. ControlNet~\cite{zhang2023adding} advances this line of research by adapting and finetuning an existing pretrained diffusion model to integrate external control cues, thus demonstrating remarkable capabilities in controllable generation. Similarly, T2I-Adapter~\cite{mou2023t2i} pursues increased image controllability by finetuning pretrained diffusion models, augmenting their responsiveness to conditioning signals. Adhering to the experimental settings outlined in \cite{zhang2023adding,mou2023t2i}, we utilize OFT for finetuning pretrained diffusion models, consistently achieving improved controllability, with the added advantage of requiring fewer labeled examples and a reduced number of parameters for finetuning. Crucially, OFT does not incur any extra computational cost during inference at test time.

\textbf{Model finetuning}. Finetuning large-scale pretrained models on a variety of downstream tasks has become widespread~\cite{devlin2018bert,brown2020language,he2022masked}. Within this context, parameter-efficient adaptation methods (\emph{e.g.}, \cite{hulora2022,houlsby2019parameter,pfeiffer2020adapterfusion}) are a focal area, especially in natural language processing. OFT bears closest resemblance to LoRA~\cite{hulora2022}, which postulates that updates to the weight matrices during finetuning can be adequately captured by a low-rank additive parameterization. Diverging from this, OFT employs a multiplicative update strategy, leveraging orthogonal transformations shared across layers, to modify neuron weights. This approach mathematically maintains the pairwise angular relationships among neurons within the same layer, leading to enhanced stability in model behavior.

\section{Orthogonal Finetuning}

\subsection{Why Does Orthogonal Transformation Make Sense?}

To motivate the adoption of orthogonal transformations for finetuning text-to-image diffusion models, we first break the problem into two parts: (1) the rationale for finetuning neuron directions (that is, their angles), and (2) the justification for choosing orthogonal transformations as the means to adjust these angles.

Addressing the first question, we build upon empirical findings from prior work~\cite{liu2018decoupled,chen2020angular}, which suggest that differences in angular features effectively encapsulate the semantic gap. SphereNet~\cite{liu2017deep} further establishes that constraining all neurons to a unit hypersphere during neural network training does not diminish capacity and may even enhance generalization, highlighting that neuron directions alone can encode essential data properties. To more clearly illustrate the significance of neuron angularity, we carry out a controlled experiment, depicted in Figure~\ref{fig:toy_ae}, where a typical convolutional autoencoder is trained on flower images. During training, we rely on the standard inner product to generate feature maps ($z$ represents the result of the convolution kernel $\bm{w}$ applied to the input $\bm{x}$ within a sliding window). At test time, we contrast three distinct approaches for computing feature maps: (a) reusing the inner product from training, (b) preserving only magnitude information, and (c) retaining solely angular (directional) information. The findings in Figure~\ref{fig:toy_ae} reveal that relying on neuron angular information nearly restores the original images, whereas the magnitude does not convey meaningful content. Notably, cosine activation (c) is excluded during training—only inner product-based learning occurs. These outcomes indicate that neuron angles, or directions, carry the bulk of the semantic information in the inputs, implying that fine-tuning neuron orientations is likely a highly effective strategy for manipulating semantic characteristics in images.

Regarding the second question, a straightforward approach to adjust neuron direction is to apply a uniform rotation or reflection to all neurons within a layer, naturally introducing an orthogonal transformation. While more nuanced angular manipulations that rotate each neuron by different amounts could, in principle, be adopted for increased expressiveness, our experiments show that orthogonal transformations offer an optimal compromise between adaptability and structural integrity. Additionally, previous studies such as~\cite{liu2021orthogonal} indicate that orthogonal transformations possess sufficient expressive power for neural network learning. To further substantiate our claim, we conduct a regularization experiment utilizing the ControlNet~\cite{zhang2023adding} framework, with outcomes presented in Figure~\ref{fig:ortho}. The data indicate that standard OFT, when employing the orthogonality constraint, maintains consistent performance and attains precise control post-training (at epoch 20). In contrast, eliminating this constraint results in the inability to generate realistic images or exert effective control. These results underscore the critical function of orthogonality in the success of OFT.

Traditional finetuning typically relies on gradient descent with a reduced learning rate to adjust the parameters of a model, or selective layers thereof. Employing a low learning rate inherently biases the optimization toward minimal departure from the pretrained state; thus, standard finetuning can be interpreted as attempting to learn model parameters while implicitly constraining $\thickmuskip=2mu \medmuskip=2mu \| \bm{M} - \bm{M}^0\|$, where $\bm{M}$ refers to the adapted model parameters and $\bm{M}^0$ to the original pretrained parameters. While this latent regularization is intuitive, it may nonetheless allow excessive degrees of freedom, especially for the finetuning of large-scale models. To make this restriction more explicit, LoRA introduces a low-rank limitation on the weight updates, namely, $\thickmuskip=2mu \medmuskip=2mu \text{rank}(\bm{M} - \bm{M}^0)=r'$ for some small $r'$. In contrast, OFT imposes a different condition: it maintains pairwise neuron similarity by enforcing $\thickmuskip=2mu \medmuskip=2mu \| \text{HE}(\bm{M})-\text{HE}(\bm{M}^0) \|=0$, where $\text{HE}(\cdot)$ is the hyperspherical energy corresponding to a weight matrix.

Consider as an example a fully connected layer represented as $\thickmuskip=2mu \medmuskip=2mu \bm{W}=\{\bm{w}_1,\cdots,\bm{w}_n\}\in\mathbb{R}^{d\times n}$, with each column $\bm{w}_i\in\mathbb{R}^{d}$ denoting an individual neuron ($\bm{W}^0$ being the pretrained weights). The output $\thickmuskip=2mu \medmuskip=2mu \bm{z}\in\mathbb{R}^n$ from this layer is computed as $\thickmuskip=2mu \medmuskip=2mu \bm{z}=\bm{W}^\top\bm{x}$, where $\thickmuskip=2mu \medmuskip=2mu \bm{x}\in\mathbb{R}^d$ is the input feature vector. From this perspective, OFT seeks to minimize the discrepancy in hyperspherical energy between the adapted and original models:

\begin{equation}\label{eq:oft_principle}
\footnotesize
\min_{\bm{W}} \left\| \text{HE}(\bm{W})-\text{HE}(\bm{W}^0)\right\|~~\Leftrightarrow~~\min_{\bm{W}} \bigg{\|} \sum_{i\neq j} \|\hat{\bm{w}}_i-\hat{\bm{w}}_j\|^{-1}-\sum_{i\neq j} \|\hat{\bm{w}}^0_i-\hat{\bm{w}}^0_j\|^{-1}\bigg{\|}
\end{equation}

Here, $\thickmuskip=2mu \medmuskip=2mu \hat{\bm{w}}_i:=\bm{w}_i/\|\bm{w}_i\|$ denotes the normalized form of neuron $i$, and hyperspherical energy for layer $\bm{W}$ is calculated as $\thickmuskip=2mu \medmuskip=2mu \text{HE}(\bm{W}):= \sum_{i\neq j} \|\hat{\bm{w}}_i-\hat{\bm{w}}_j\|^{-1}$. It is straightforward to note that the optimal solution for Eq.~\eqref{eq:oft_principle} yields a minimum value of zero. This is achieved if $\bm{W}$ is related to $\bm{W}^0$ via an orthogonal transformation, that is, $\thickmuskip=2mu \medmuskip=2mu \bm{W}=\bm{R}\bm{W}^0$, where $\thickmuskip=2mu \medmuskip=2mu \bm{R}\in\mathbb{R}^{d\times d}$ is orthogonal (with determinant $1$ or $-1$ representing rotation or reflection, respectively). The core premise of OFT is thus to finetune the network by learning layer-specific orthogonal transformations applied to all

\begin{equation}\label{eq:oft_general}
    \footnotesize
    \bm{z}=\bm{W}^\top\bm{x}=(\bm{R}\cdot\bm{W}^0)^\top\bm{x},~~~\text{s.t.}~\bm{R}^\top\bm{R}=\bm{R}\bm{R}^\top=\bm{I}
\end{equation}

Here, $\bm{W}$ stands for the weight matrix obtained after OFT finetuning, while $\bm{I}$ represents the identity matrix. An illustration of the OFT approach is provided in Figure~\ref{fig:block}. Analogous to LoRA’s use of zero initialization, it is essential that OFT allows finetuning to commence precisely from the pretrained parameters $\bm{W}^0$. To realize this, we set the initial value of the orthogonal matrix $\bm{R}$ to be the identity matrix, ensuring the finetuned model parameters are initially identical to those of the pretrained model. The orthogonality constraint on $\bm{R}$ can be maintained by leveraging differentiable orthogonalization methods as described in \cite{liu2021orthogonal,lezcano2019cheap}. Later in this section, we detail approaches to enforcing orthogonality in a computationally practical manner.

\subsection{Efficient Orthogonal Parameterization}\label{sect:eff_ortho}

Traditional approaches for orthogonalization, such as the Gram-Schmidt process, although differentiable, are frequently computationally intensive and impractical for large models~\cite{liu2021orthogonal}. To enhance efficiency, we make use of the Cayley parameterization to construct the orthogonal matrix. Concretely, the orthogonal matrix can be written as $\thickmuskip=2mu \medmuskip=2mu \bm{R}=(\bm{I}+\bm{Q})(I-\bm{Q})^{-1}$, where $\bm{Q}$ is a skew-symmetric matrix such that $\thickmuskip=2mu \medmuskip=2mu \bm{Q}=-\bm{Q}^\top$. While this Cayley parameterization notably accelerates computation, it does restrict the resulting orthogonal matrices to have a determinant of $1$, placing them within the special orthogonal group. However, empirical results suggest that this restriction does not compromise practical performance. Despite the adoption of the Cayley transform, parameterizing the full orthogonal matrix $\bm{R}$ remains inefficient, particularly for large $d$. To mitigate this, we introduce a block-diagonal structure for $\bm{R}$ composed of $r$ blocks, yielding the following representation:

\begin{equation}
    \footnotesize
    \bm{R}=\text{diag}(\bm{R}_1,\bm{R}_2,\cdots,\bm{R}_r)=\begin{bmatrix}
    \bm{R}_1\in \text{O}(\frac{d}{r}) &  &\\
    &\ddots & \\
    & & \bm{R}_r\in \text{O}(\frac{d}{r})
    \end{bmatrix}\in \text{O}(d)
\end{equation}

Here, $\text{O}(d)$ represents the orthogonal group of degree $d$, with $\thickmuskip=2mu \medmuskip=2mu \bm{R}\in\mathbb{R}^{d\times d}$ and each $\thickmuskip=2mu \medmuskip=2mu \bm{R}_i\in\mathbb{R}^{d/r\times d/r},\forall i$, being orthogonal matrices. In the special case where $\thickmuskip=2mu \medmuskip=2mu r=1$, the block-diagonal orthogonal structure degenerates to the standard unconstrained orthogonal matrix. For an orthogonal matrix of size $\thickmuskip=2mu \medmuskip=2mu d\times d$, there are $\thickmuskip=2mu \medmuskip=2mu d(d-1)/2$ free parameters, yielding a parameter complexity of $\mathcal{O}(d^2)$. In contrast, if the orthogonal matrix is $r$-block diagonal, the parameter count reduces to $\thickmuskip=2mu \medmuskip=2mu d(d/r-1)/2$, so the complexity is $\thickmuskip=2mu \medmuskip=2mu \mathcal{O}(d^2/r)$. An additional reduction is possible by tying the blocks together, that is, by setting $\thickmuskip=2mu \medmuskip=2mu\bm{R}_i=\bm{R}_j$ for all $i\neq j$, which drops the parameter complexity to $\mathcal{O}(d^2/r^2)$. Crucially, all these approaches preserve the orthogonality of the overall matrix $\bm{R}$, ensuring that hyperspherical energy is maintained.

We next compare the parameter efficiency of OFT to that of LoRA. In LoRA, for a rank $r'$ decomposition, the trainable parameter count is $\thickmuskip=2mu \medmuskip=2mu r'(d+n)$. If both $r$ and $r'$ are scaled with the neuron dimension $d$ (for instance, $\thickmuskip=2mu \medmuskip=2mu r=r'=\alpha d$, with $\thickmuskip=2mu \medmuskip=2mu 0<\alpha\leq 1$ a fixed coefficient), then LoRA's total parameter complexity becomes $\mathcal{O}(d^2+dn)$, whereas for OFT it reduces to $\mathcal{O}(d)$. This efficiency gap is illustrated by a practical case: with a weight matrix sized $\thickmuskip=2mu \medmuskip=2mu 128\times 128$, LoRA with $\thickmuskip=2mu \medmuskip=2mu r'=8$ results in $2,048$ learnable parameters, while OFT with $\thickmuskip=2mu \medmuskip=2mu r=8$ (without block sharing) requires just $960$ trainable parameters.

\subsection{Constrained Orthogonal Finetuning}

The expressivity of the original OFT method can be further narrowed by enforcing that the fine-tuned model stays close to the initial pre-trained solution. In particular, COFT employs the following inference step:

\begin{equation}\label{eq:coft_general}
    \footnotesize
    \bm{z}=\bm{W}^\top\bm{x}=(\bm{R}\cdot\bm{W}^0)^\top\bm{x},~~~\text{s.t.}~\bm{R}^\top\bm{R}=\bm{R}\bm{R}^\top=\bm{I},~\norm{\bm{R}-\bm{I}}\leq \epsilon
\end{equation}

Here, the orthogonality constraint is enforced on $\bm{R}$, along with a bound on its distance from the identity matrix given by the $\epsilon$-deviation constraint. The Cayley parameterization, as described in Section~\ref{sect:eff_ortho}, provides a suitable means of imposing orthogonality. However, integrating the $\epsilon$-deviation directly into a Cayley-parameterized matrix is non-trivial. To further understand the Cayley transform's properties, consider the Neumann series expansion: with $\thickmuskip=2mu \medmuskip=2mu \bm{R}=(\bm{I}+\

Given that the series converges in the operator norm, we are justified in incorporating the constraint $\thickmuskip=2mu \medmuskip=2mu\|\bm{R}-\bm{I}\|\leq\epsilon$ directly into the Cayley transform. This leads to a corresponding restriction on $\bm{Q}$ of the form $\thickmuskip=2mu \medmuskip=2mu \|\bm{Q}-\bm{0}\|\leq\epsilon'$, with $\bm{0}$ representing the zero matrix and $\epsilon'$ serving as a distinct error tolerance parameter. This revised constraint involving $\bm{Q}$ lends itself well to enforcement via projected gradient descent. To ensure that the orthogonal matrix $\bm{R}$ is initialized to the identity, we begin with $\bm{Q}$ initialized as the zero matrix. The COFT approach therefore imposes two direct constraints: a restriction on hyperspherical energy variation and a regulated deviation from the pretrained parameters. While this latter requirement is commonly an implicit part of standard finetuning routines, COFT makes it explicit. Although OFT has demonstrated excellent performance, our empirical observations reveal that COFT offers enhanced stability during finetuning, attributable to the explicit constraint on parameter drift. Figure~\ref{fig:eps_coft} illustrates the sensitivity of COFT to the choice of $\epsilon$. This parameter modulates the permissible adaptation: increasing $\epsilon$ allows COFT to recover behavior similar to OFT, whereas decreasing $\epsilon$ causes COFT to increasingly approximate the original pretrained text-to-image diffusion model.

\subsection{Re-scaled Orthogonal Finetuning}

We introduce an intuitive enhancement to the original OFT by incorporating a learnable scaling factor for each neuron. The rationale is that neuron rescaling leaves the hyperspherical energy unchanged, since the scaling effect is removed during normalization. Formally, the modified forward operation is given by $\thickmuskip=2mu \medmuskip=2mu\bm{z}=(\bm{R}\bm{W}^0\bm{D})^\top\bm{x}$,\footnote{\textbf{Errata}: The camera-ready NeurIPS paper incorrectly stated the forward computation for re-scaled OFT as $\thickmuskip=2mu \medmuskip=2mu\bm{z}=(\bm{D}\bm{R}\bm{W}^0)^\top\bm{x}$. The implementation was correct, so results in Appendix~\ref{app:rescaled} remain valid.} where $\thickmuskip=2mu \medmuskip=2mu \bm{D}=\text{diag}(s_1,\cdots,s_n)\in\mathbb{R}^{n\times n}$ is a trainable diagonal matrix with positive entries $s_1,\ldots,s_n$. Unlike the original OFT, as defined in Eq.~\eqref{eq:oft_general}, which updates only $\bm{R}$, the re-scaled variant simultaneously learns both the orthogonal matrix $\bm{R}$ and the scaling matrix $\bm{D}$. This adjustment expands OFT’s expressiveness with a minimal parameter overhead. In our main experiments, we retain the original OFT setup to highlight the benefits of orthogonal transformations alone, though empirical results (see Appendix~\ref{app:rescaled}) generally favor the re-scaled version.

\section{Intriguing Insights and Discussions}

\textbf{OFT’s architectural generality.} In principle, OFT can be adapted to any neural network architecture. For Transformers, LoRA customarily targets the attention weight matrices~\cite{hulora2022}; to provide a direct comparison, our experiments restrict OFT application to the attention weights as well. Additionally, OFT extends naturally to convolutional layers, in part due to the interpretability of the block-diagonal form of $\bm{R}$—a property not shared by LoRA. Specifically, matching the number of blocks to the number of input channels leads each block to affect a single neuron channel, akin to implementing learnable depth-wise convolution kernels~\cite{chollet2017xception}. If instead all blocks of $\bm{R}$ are identical, OFT employs the same orthogonal transformation across all channels. Further exploratory experiments on OFT-based finetuning of convolutional layers are included in Appendix~\ref{app:convolution}.

\textbf{Connection to LoRA}. LoRA incorporates a low-rank matrix to ensure that modifications to the pretrained weight matrix are minimal, thereby preserving its original information. In contrast, OFT constrains the transformation applied to the pretrained weights to remain orthogonal and full-rank, which also helps to protect the integrity of the pretrained information. The forward computation in OFT can be reformulated as $\thickmuskip=2mu \medmuskip=2mu \bm{z}=(\bm{R}\bm{W}^0)^\top\bm{x}=(\bm{W}^0+(\bm{R}-\bm{I})\bm{W}^0)^\top\bm{x}$, wherein $\thickmuskip=2mu \medmuskip=2mu (\bm{R}-\bm{I})\bm{W}^0$ serves a role comparable to LoRA's low-rank adjustment to the weights. Given that $\bm{W}^0$ generally possesses full rank, OFT effectively induces a low-rank update if $\thickmuskip=2mu \medmuskip=2mu \bm{R}-\bm{I}$ itself is low-rank. Like the rank parameter $r'$ in LoRA, OFT introduces a block-diagonal size $r$ that controls the number of trainable parameters. Notably, the two approaches exemplify fundamentally different parameter-efficient strategies: LoRA leverages a low-rank structure to reduce trainable variables, while OFT benefits from sparsity, namely block-diagonal orthogonal matrices, to achieve a similar goal.

\textbf{Why OFT converges faster}. On the one hand, Figure~\ref{fig:toy_ae} suggests that the optimal way to change model semantics is by rotating neuron directions—precisely what OFT is structured to accomplish. On the other hand, since OFT can be interpreted as optimizing neurons constrained to a smooth hypersphere manifold, it results in a more favorable optimization landscape. This claim is also supported by the empirical results reported in \cite{liu2021orthogonal}.

\textbf{Why not minimize hyperspherical energy}. Unlike the approach of \cite{liu2021orthogonal}, the objective here is not to minimize hyperspherical energy. In classification tasks, it is preferable for neurons to have little redundancy, but arranging neurons to be evenly distributed across the hypersphere—achieving minimal hyperspherical energy—does not align with the aims of finetuning, as this process could disrupt critical information learned during pretraining.

\textbf{Balancing flexibility and regularity in finetuning}. Our study highlights a fundamental tension between flexibility and regularization in model finetuning. While conventional finetuning is highly adaptable, it often compromises stability and can lead to model degeneration. In contrast, OFT, though conceptually straightforward, effectively mediates this balance by maintaining the angular relationships between neurons. Additionally, employing a block-diagonal structure serves as a form of stronger regularization applied to the orthogonal transformation.

\textbf{No extra computational cost during inference}. In contrast to ControlNet, the OFT approach does not impose any additional computational burden during the inference phase. During inference, the learned orthogonal matrix $\bm{R}$ can be directly combined with the pretrained weights $\bm{W}^0$, resulting in the composite weight matrix $\thickmuskip=2mu \medmuskip=2mu \bm{W} = \bm{R} \bm{W}^0$. Consequently, inference throughput remains unchanged relative to that of the original pretrained model.

\section{Experiments and Results}

\textbf{General experimental protocol}. Our experiments utilize Stable Diffusion v1.5~\cite{rombach2022high} as the foundational text-to-image model. For objective assessment, outputs are randomly sampled across competing methods. Subject-driven generation follows protocols from DreamBooth~\cite{ruiz2023dreambooth}, while approaches for controllable image synthesis adhere to ControlNet~\cite{zhang2023adding} and T2I-Adapter~\cite{mou2023t2i}. To equitably benchmark OFT or COFT against LoRA, all approaches are applied to identical network layers. Additional implementation specifics and results are available in Appendix~\ref{app:settings}.

\subsection{Subject-driven Generation}

\textbf{Experimental settings}. DreamBooth~\cite{ruiz2023dreambooth} and LoRA~\cite{hulora2022} serve as primary baselines in our comparison. Each method adopts the loss formulation presented in DreamBooth. Experimental configurations for DreamBooth and LoRA closely mirror their respective original works, leveraging optimal hyperparameter choices. Supplementary findings can be found in Appendix~\ref{app:settings},~\ref{app:coft_oft},~\ref{app:more_qualitative}, and~\ref{app:failure}.

\textbf{Finetuning stability and convergence}. We begin by examining the robustness and speed of convergence during finetuning for DreamBooth, LoRA, OFT, and COFT. The corresponding findings are illustrated in Figure~\ref{fig:teaser} and Figure~\ref{fig:db_convergence}. Both COFT and OFT display stable finetuning behavior when applied to diffusion models. After 400 finetuning steps, effective subject control is observed for both DreamBooth and the OFT approaches, whereas LoRA struggles with subject identity retention. By 2000 steps, DreamBooth begins to exhibit mode collapse, generating low-quality outputs, and LoRA inaccurately renders the subject’s clothing (producing yellow fur instead of a yellow shirt). In contrast, OFT and COFT continue to deliver reliable and consistent image control throughout the training process. These outcomes underscore the capacity of our OFT methodology to rapidly reach stable solutions in subject-driven generation tasks. An additional advantage lies in the minimal sensitivity of OFT and COFT to the total number of finetuning iterations, greatly simplifying the process of determining a suitable iteration count across different subjects. As a result, a generously large iteration limit can be selected by default, without the need for extensive hyperparameter adjustment. With an appropriately chosen $\epsilon$ for COFT, both the learning rate and the number of finetuning iterations require little manual tuning.

\textbf{Quantitative comparison}. Adopting the protocol from \cite{ruiz2023dreambooth}, we quantitatively assess the models based on subject fidelity (using DINO~\cite{caron2021emerging} and CLIP-I~\cite{radford2021learning}), alignment with text prompts (CLIP-T~\cite{radford2021learning}), and sample diversity (LPIPS~\cite{zhang2018unreasonable}). In this context, CLIP-I represents the mean pairwise cosine similarity between the CLIP embeddings of generated and real images, while DINO mirrors this metric using ViT S/16 DINO features. For CLIP-T, the average cosine similarity is measured between text prompt embeddings and those of generated images. Additionally, LPIPS evaluates the mean cosine similarity between multiple generated samples of the same subject given identical prompts. Table~\ref{table:dreambooth_final} demonstrates that COFT and OFT significantly surpass DreamBooth and LoRA on the DINO and CLIP-I metrics and yield marginally superior or comparable results for prompt fidelity and diversity. For rigorous comparison, each method is run for 30 random seeds during finetuning of the same base model, reducing stochastic variability. Overall, our experiments confirm that OFT consistently achieves superior convergence stability and delivers enhanced final performance compared to alternative methods.

\textbf{Qualitative comparison}. To provide clearer insight into the advantages offered by OFT, we present a selection of randomly chosen subject-driven generation samples in Figure~\ref{fig:dreambooth_final}. In order to maintain fairness, each sample is produced by the same finetuned model across all methods; importantly, we refrain from custom-tailoring text prompts to enhance any individual result. For each approach, we pick the model checkpoint that attains the highest validation CLIP score. As illustrated in Figure~\ref{fig:dreambooth_final}, both OFT and COFT consistently excel at retaining the subject's semantic identity, whereas LoRA frequently struggles to maintain subject fidelity (for example, in the bowl scenario, LoRA completely fails to preserve the subject's identity). Furthermore, OFT and COFT provide more reliable responsiveness to text prompts, while DreamBooth, although it effectively preserves subject identity, often falls short in generating content faithful to the prompt (see, for instance, the first row for the bowl example). Overall, this qualitative assessment indicates that OFT yields improved balance between controllability and subject fidelity. Additionally, the OFT framework shows robustness to iteration count: its performance remains strong even with an increased number of iterations, a property not shared by DreamBooth or LoRA. Supplementary qualitative results can be found in Appendix~\ref{app:more_qualitative}. In addition, a human study reported in Appendix~\ref{app:human} further supports the effectiveness of OFT.

\subsection{Controllable Generation}

\textbf{Settings}. For our baselines, we consider ControlNet~\cite{zhang2023adding}, T2I-Adapter~\cite{mou2023t2i}, and LoRA~\cite{hulora2022}. The principal evaluation comprises three complex controllable generation benchmarks: Canny edge to image~(C2I) on COCO~\cite{lin2014microsoft}, segmentation map to image~(S2I) on ADE20K~\cite{zhou2017scene}, and landmark to face~(L2F) using CelebA-HQ~\cite{karras2018progressive,xia2021tedigan}. Each approach is used to finetune Stable Diffusion~(SD) v1.5 for 20 epochs on each of these datasets. Expanded experimental outcomes are provided in Appendix~\ref{app:more_qualitative},~\ref{app:more_control_task}, and~\ref{app:failure}.

\textbf{Convergence}. We investigate how rapidly ControlNet, T2I-Adapter, LoRA, and COFT converge within the context of the L2F task, employing both quantitative and qualitative assessments. For our primary metric, we calculate the average $\ell_2$ distance separating the ground truth face landmarks and those predicted by each model. Figure~\ref{fig:face_convergence} presents a comparison of landmark errors as training proceeds over varying numbers of epochs for each fine-tuned model. Our analysis reveals that COFT and OFT both display notably accelerated convergence relative to baselines. Specifically, LoRA requires 20 epochs to reach the same performance level that OFT attains by epoch 8. It is important to emphasize that OFT and COFT utilize a number of trainable parameters similar to LoRA (and substantially fewer than ControlNet), yet converge much more rapidly than the current alternatives. The rapid progress made by OFT is further substantiated by the examples shown in Figure~\ref{fig:teaser}: here, OFT demonstrates a high degree of data efficiency, successfully converging even when trained with only 5\% of the ADE20K dataset, outpacing both ControlNet and LoRA. For qualitative analysis, we limit our comparison to OFT, COFT, and ControlNet since the latter provides the most comparable landmark accuracy. As illustrated in Figure~\ref{fig:face_convergence_qual}, OFT and COFT consistently exhibit stable convergence, with generated facial poses smoothly aligning with the provided control landmarks over training. In contrast, ControlNet shows a sudden improvement in controllability after epoch 8, mirroring the abrupt convergence described in \cite{zhang2023adding}. The steady and smooth training behavior observed in OFT greatly enhances its practical usability, making its learning process more reliable and facilitating the selection of effective hyperparameters.

\textbf{Quantitative comparison}. To assess the effectiveness of controllable generation, we propose a control consistency metric, which operates by extracting the control signal from the output image and comparing it to the original input control. Specifically, for the C2I task, the evaluation employs IoU and F1 scores. The S2I task is assessed via mean IoU, mean accuracy, and overall accuracy. For the L2F task, we utilize the mean $\ell_2$ distance between the actual control landmarks and those predicted. Further explanations on these consistency metrics are included in Appendix~\ref{app:settings}. All competing approaches are benchmarked using their optimal hyperparameter configurations. The outcomes, detailed in Table~\ref{table:control_gen}, indicate that OFT and COFT consistently provide more precise and reliable control than alternative models. Notably, adapter-based methods such as T2I-Adapter and ControlNet demonstrate slower convergence and achieve inferior final performance. LoRA surpasses ControlNet in S2I tasks but underperforms in both C2I and L2F scenarios. Broadly, even with careful hyperparameter optimization, LoRA’s maximum achievable performance remains modest. In contrast, our empirical observations suggest that the OFT framework’s performance continues to improve as the count of trainable parameters increases, showing no signs of saturation. Finally, we highlight that our quantitative methodology for assessing controllable generation is a novel contribution, enabling precise evaluation of the control capabilities of finetuned models, within the bounds set by the pre-existing segmentation or detection models.

\textbf{Qualitative comparison}. We conduct a qualitative analysis of OFT and COFT against leading approaches such as ControlNet, T2I-Adapter, and LoRA. As depicted by the random samples in Figure~\ref{fig:control_final}, both OFT and COFT consistently deliver images with impressive realism and high fidelity, alongside precise control capabilities. In the S2I task, LoRA is unable to reliably synthesize images that align with the provided segmentation maps, whereas ControlNet, OFT, and COFT exhibit strong control over image generation. Notably, OFT and COFT surpass ControlNet by producing more detailed and geometrically consistent images, despite utilizing significantly fewer parameters. Within the C2I scenario, OFT and COFT demonstrate the capacity to synthesize semantically meaningful images from rudimentary Canny edge maps, a feat where T2I-Adapter and LoRA are less effective. For the L2F task, our approach achieves the highest degree of accuracy in facial pose control, even for complex head positions. Across all three control scenarios, OFT and COFT consistently outperform baseline models in visual quality, underscoring the robust controllability of the OFT framework. To further strengthen our qualitative assessment, we present additional visual examples for all three control tasks in Appendix~\ref{app:control_exp}. Moreover, Appendix~\ref{app:more_control_task} showcases the versatility of OFT across a wider range of control tasks, such as dense pose to human body, sketch to image, and depth to image generation.

\section{Concluding Remarks and Open Problems}

Inspired by the insight that angular relationships among neurons play a pivotal role in shaping visual semantics, we introduce a straightforward yet robust finetuning technique—orthogonal finetuning (OFT)—designed to modulate text-to-image diffusion models. Our approach is particularly applied to two scenarios: subject-driven generation and controllable generation. In comparison to prior methods, OFT achieves enhanced control and finetuning stability, while utilizing a smaller set of finetuning parameters. Notably, OFT avoids adding inference complexity, making the model both efficient and suitable for deployment.

Several compelling questions arise from the use of OFT. To ensure orthogonality, OFT relies on Cayley parametrization, which necessitates computing a matrix inverse and thus poses a scalability constraint. While we mitigate this issue through block diagonal parametrization, developing differentiable and faster techniques for matrix inversion remains an unresolved challenge. Furthermore, OFT presents interesting prospects for compositionality: the orthogonal matrices obtained from distinct OFT finetuning tasks can be combined through multiplication, and the result remains orthogonal. Investigating whether these composed orthogonal matrices retain knowledge from all associated downstream tasks is a promising research avenue. Lastly, although the parameter efficiency of OFT is mainly due to the block diagonal design, this also introduces some inherent biases and limits adaptability. Exploring more advanced and less biased methods for improving parameter efficiency stands out as an important unresolved question.